\chapter{Introduction}\label{ch:introduction}\pagenumbering{arabic}

Creating backups is important.
With the possibility of a catastrophic error such as a hardware failure happening
at any time keeping data save is a known challenge with a possible solution
being periodic backups.

Since periodic backups often contain large amounts of data which did not change
since the previous backup storage space is wasted on storing redundant data.
Many existing solutions are able to remove redundant data from individual backups
in a process called `deduplication'.
The storage space savings created by such an ability can be increased by performing
it across the maximum number of backups possible.

This in turn is limited by many current implementations placing restrictions on
the ability to perform deduplication across backups originating from different
sources.
Should for example all workstation machines existing in an office be backed up
to a central server there can be restrictions such as deduplication only being
possible between backups of the same machine or only one machine being able to
create a backup at a time.

To prevent having to implement such restrictions while also ensuring backup
data remains consistent Gilbert et al.\ developed an algorithm called
`Lock-Free Deduplication'~\cite{Duplicacy}.
Its greatest difference in comparison to existing approaches is the ability of
all operations to be performed in parallel, even by different machines accessing
a shared backup storage without sacrificing deduplication.

While the algorithm is implemented in a program called `Duplicacy' its usefulness
is limited by the available interfaces and licenses.
With Duplicacy being intended for use by end users it is difficult to reuse the
algorithm in other programs or approaches besides managing backup data.
Reuse is further complicated by its license\footnote{\url{https://github.com/gilbertchen/duplicacy/blob/master/LICENSE.md}}
requiring payment for commercial use for each machine running it.

This thesis concerns itself with the question on how to support alternative
implementations of the Lock-Free Deduplication algorithm.
First, there will be an examination of existing approaches to creating backups
and the details of the algorithm in \Cref{ch:background}.
Next, problems with applying it to real-world conditions are identified in
\Cref{ch:design} and solutions provided.
After that a library called `Vinculum' which is the main result of this thesis
will be presented in \Cref{ch:implementation}.
Its performance will be evaluated in \Cref{ch:evaluation} by comparing it with
a selection of similar backup solutions followed by a discussion of the results
in \Cref{ch:discussion}.
Finally, \Cref{ch:conclusion} concludes the thesis.
